\documentclass[11pt]{article}
\usepackage[utf8]{inputenc}
\usepackage{amsmath,amssymb}
\usepackage[margin=1in]{geometry}
\usepackage{hyperref}
\emergencystretch=3em

\title{The multiplicity-$m$ theorem for an arbitrary coordinate order}
\author{Claude, with Daniel Murfet}
\date{2026-09-07}

\begin{document}
\maketitle
\sloppy

\begin{abstract}
Relabelling the coordinates of the box along an equivalence $\tau:\mathrm{Fin}\,D'\simeq\mathrm{Fin}\,D$
leaves the chart integral invariant when $\xi,\eta$ are transported along the induced measurable
equivalence. Sorting the exponents puts the minimisers last, so unit~80 yields the general form of the
leading term of \texttt{thm:TaylorTree}: for any exponents on $\mathrm{Fin}\,D$ whose minimal candidate
exponent $p$ is attained $M+2\ge2$ times, and jointly continuous $\xi,\eta$ with $\eta>0$ on the minimal
face, $\int_{(0,b]^D}w^H\eta\,e^{-\beta nw^{2K}+\beta\sqrt n\,w^K\xi}\sim C\,n^{-p/2}\log^{M+1}n$ with
$C>0$. Zero \texttt{sorry}/\texttt{axiom}.
\end{abstract}

\section{The things formalised}
{\small
\begin{verbatim}
noncomputable def relabel {D D' : ℕ} (τ : Fin D' ≃ Fin D) : (Fin D' → ℝ) ≃ᵐ (Fin D → ℝ)
theorem boxIntegralGen_relabel (β b c : ℝ) {D D' : ℕ} (τ : Fin D' ≃ Fin D) (K H : Fin D → ℕ)
    (ξ η : (Fin D → ℝ) → ℝ) :
    boxIntegralGen β b c (K ∘ τ) (H ∘ τ) (ξ ∘ relabel τ) (η ∘ relabel τ)
      = boxIntegralGen β b c K H ξ η
theorem boxIntegralGen_isEquivalent_general (β b : ℝ) (hβ : 0 < β) (hb : 0 < b) {D : ℕ}
    (K H : Fin D → ℕ) (hK : ∀ i, 0 < K i) (p : ℝ) (hmin : ∀ i, p ≤ finExp K H i) (M : ℕ)
    (hM : (Finset.univ.filter fun i => finExp K H i = p).card = M + 2)
    (ξ η : (Fin D → ℝ) → ℝ) (hξc : Continuous ξ) (hηc : Continuous η)
    (hηpos : ∀ w : Fin D → ℝ, (∀ i, finExp K H i = p → w i = 0) →
      (∀ i, finExp K H i ≠ p → 0 < w i ∧ w i ≤ b) → 0 < η w) :
    ∃ C : ℝ, 0 < C ∧ (fun n : ℝ => boxIntegralGen β b (Real.sqrt n) K H ξ η) ~[atTop]
      fun n : ℝ => C * (n ^ (-(p / 2)) * Real.log n ^ (M + 1))
\end{verbatim}
}

\section{Proof strategy}
Relabelling invariance is \texttt{measurePreserving\_piCongrLeft} applied to
\texttt{MeasurePreserving.integral\_comp'}, with the monomials reindexed by
\texttt{Fintype.prod\_equiv}. For the general theorem, $\sigma=\mathrm{sortPerm}$ makes the exponents
antitone; the minimisers form the final segment $\mathrm{Ici}\,i_0$ of cardinality $M+2$, so
$i_0=D-(M+2)=:d'$, and $\tau=\mathrm{finCongr}\circ\sigma$ relabels $\mathrm{Fin}(d'+(M+2))$. The
exponent hypotheses of unit~80 follow from the segment description; the face condition transports
through \texttt{relabel\_apply\_apply} and \texttt{Fin.addCases} on $\tau^{-1}i$.

\section{Roadblocks and resolutions}
\begin{itemize}
\item \texttt{MeasurableEquiv.piCongrLeft\_apply\_apply} needs the family given explicitly
(\texttt{(β := fun \_ => ℝ)}) when the equivalence is wrapped in a definition.
\item \texttt{Finset.card\_equiv} maps the left set into the right one: with the sorted set on the
right the equivalence is $\sigma^{-1}$, not $\sigma$.
\item \texttt{Fin.addCases} does not generalise hypotheses mentioning the eliminated index;
\texttt{revert} them first.
\end{itemize}

\section{Outcome}
The leading term of \texttt{thm:TaylorTree} is formalised for general continuous $\xi,\eta$ on the
genuine chart box in any coordinate order and any multiplicity ($m=1$ by unit~72, $m\ge2$ here).
Next: consult Astra on the remainder and expansion programme (targets~6--8).
\end{document}
